\chapter[Synthetic Benchmarking]{Synthetic Benchmarking: Triage Before
You Buy GPU Time}
\label{ch:synthetic}

An architecture evaluation on real data costs GPU-days and returns one
noisy bit. When forensics has recovered the shape of the data-generating
process (Part~II), a better instrument exists: \textbf{build a synthetic
world with that shape, plant mechanisms you control, compute the optimal
score in closed form, and triage architectures against the known ceiling
--- at laptop cost}. This chapter presents the method and the generator we
built (\code{scripts/synth\_bench.py}).

\section{Why synthetic tests usually lie, and how to make them honest}

The standard objection is sound: toy benchmarks reward toy strengths.
Synthetic triage earns trust only under four disciplines, each fixing a
standard failure:

\begin{enumerate}
\item \textbf{Match the SNR.} The classic sin is testing at $R^2 = 0.5$
  where clever architectures shine, then deploying at $R^2 = 0.01$ where
  regularization dominates everything. Calibrate the planted signal so
  the \emph{optimal} score sits near the real problem's ceiling (ours:
  reference $R^2 \approx 0.03$--$0.05$ against the real $0.014$ ---
  slightly generous so verdicts resolve at small data scale, and stated
  as such).
\item \textbf{Match the protocol.} Run the synthetic world through the
  \emph{identical} production harness --- same feature pipeline, same
  frozen scalers, same walk-forward with daily online refit
  (Chapter~\ref{ch:validation}). We achieve this mechanically: the
  generator writes the canonical competition file layout, so every
  existing tool runs on it unchanged.
\item \textbf{Match the information sets.} Every candidate model receives
  identical inputs (including the lagged-target features), so a win is
  attributable to the architecture and nothing else.
\item \textbf{Fix the decision rule in advance.} Ours, written before any
  run: \emph{an architecture that cannot beat the GRU on a world designed
  for its own inductive bias does not receive real-data GPU budget.}
\end{enumerate}

\section{Designing the world: one mechanism per inductive bias}

The generator simulates the recovered construction of
Chapter~\ref{ch:targets} --- per-symbol innovation series $u$, responders
as forward SMAs (4/20/120) on two noisy venues, features as noisy
descriptions of $u$'s realized path, planted at the same indices with the
same timing signatures the atlas found. Into $u$'s conditional mean we
plant four \emph{toggleable} mechanisms, each the cleanest possible target
for one architectural bias:

\begin{center}
\small
\begin{tabularx}{\textwidth}{@{}l>{\raggedright\arraybackslash}X>{\raggedright\arraybackslash}X@{}}
\toprule
mechanism & construction & should be won by \\
\midrule
\code{rev} & intraday mean reversion:
$\E[u_t] \propto -\,\mathrm{SMA}_{20}(u)_{t-1}$ & any causal temporal
model \\
\code{xd} & cross-day: yesterday's last-120 trend $\times$ a
time-of-day gate ($+$ early, $-$ late) & TimeXer's endogenous-patch
mechanism (Chapter~\ref{ch:attention}) \\
\code{gate} & $\E[u_t] \propto \min(f_a, f_b)$ of two persistent
observable factors (one idiosyncratic, one market-wide) & iTransformer's
variate attention; also what the symbolic search found on real data \\
\code{lead} & two features observe the next-4 innovations noisily
(the \feat{37}/\feat{38} analog) & anything --- the linear-Gaussian floor \\
\bottomrule
\end{tabularx}
\end{center}

The craft is in the details that make mechanisms \emph{forecastable but
not trivial}: the gate factors follow slow AR(1) dynamics
($\rho = 0.995$) so their state persists across the 20-step target window
(otherwise the mechanism is unlearnable from time $t$); the cross-day
mechanism requires \emph{composing} two inputs (yesterday's path and the
clock) that no single feature exposes; the lead features carry calibrated
noise so their optimal use is a shrinkage the model must learn.

\section{The reference predictor: score as a fraction of knowable}

For each mechanism we derive $\E[y_t \mid \mathcal F_t]$ in closed form
(exactly for \code{rev}/\code{xd}/\code{lead}; a $\rho^k$-decay
approximation for the min-gate, whose Jensen gap we accept and disclose).
The generator writes these \emph{per-mechanism oracle columns} beside the
data. Every model's score is then reported as a \textbf{fraction of the
reference}, and --- the sharper instrument --- per-mechanism attribution
is available by construction: re-generate the world with one mechanism's
coefficient zeroed and difference the scores.

\begin{keybox}[: triage reads, decided in advance]
The bench's output is a (model $\times$ world) table read against three
patterns. \emph{TimeXer $\approx$ GRU on the full world}: its bias is not
realizable at this SNR and data volume $\to$ no GPU budget --- decision
made for a laptop-hour. \emph{TimeXer $>$ GRU on full, $=$ on
xd-zeroed}: the mechanism is real and TimeXer harvests it $\to$ the
real-data question becomes ``how much \code{xd}-like structure does the
real panel have?'', which the lag ablations already bound
(Chapter~\ref{ch:featurepatterns}). \emph{Everything far below
reference}: the harness or the calibration is broken --- fix before
concluding anything about architectures.
\end{keybox}

\section{Validating the validator}

A synthetic world is itself code that can be wrong, so it must pass its
own forensics before grading anyone. Ours ships three self-checks: the
generated responders reproduce the triangular ACF ramps at the correct
cutoffs; the in-sample reference $R^2$ decomposes across mechanisms near
the calibrated targets (measured: full $+0.044$ = rev $0.014$ + xd
$0.007$ + gate $0.007$ + lead $0.011$, up to interaction terms); and the
atlas pipeline of Chapter~\ref{ch:featforensics}, run on the synthetic
world, recovers the planted feature taxonomy --- the forensic tools and
the generator audit each other.

\section{Beyond triage: the other uses}

Once the world exists, it answers questions the real data cannot:

\begin{itemize}
\item \textbf{Procedure tuning with ground truth}: refit-lr cliffs
  (Chapter~\ref{ch:online}), window sweeps, embargo sizes --- all
  measurable against a \emph{known} optimum rather than a noisy
  empirical best.
\item \textbf{Power analysis}: how many days of validation are needed to
  resolve a $+0.001$ true difference at this SNR? Simulate; the answer
  disciplines real-data conclusions
  (Chapter~\ref{ch:diagnostics}).
\item \textbf{Failure rehearsal}: inject a regime break and watch each
  family's static/online scores respond; learn the curve signatures
  cheaply.
\item \textbf{Leakage fire drills}: deliberately leak (center a window,
  drop the embargo) and confirm your audit protocol catches it. An
  audit that has never caught a planted leak is untested equipment.
\end{itemize}

\begin{drillbox}
Write the minimal version for your own problem: (1) simulate the
construction your forensics recovered, at matched SNR; (2) plant one
mechanism per architecture you are weighing; (3) derive or approximate
the reference predictor and \emph{calibrate} until its score sits near
your real ceiling; (4) run your production harness unchanged. Budget a
day. Compare that to the GPU-weeks of undirected architecture search it
replaces --- and note that the discipline of step 3 will, by itself,
teach you more about your problem's knowable fraction than most modeling
weeks do.
\end{drillbox}
